% GENERATED FILE. Edit canonical lessons, then run npm run generate:books.
% canonical-source-hash: fnv1a64:939be64e97d57d2f
% canonical-lessons: MR-C392-kalam, MR-C392-vatap-karne, MR-C392-davalne, MR-C392-anusarne, MR-C392-angikarne

\chapter{Rules in Practice}
\label{ch:mr-rules-in-practice}
\hlchaptermodality{392}

\begin{chapteropening}
\textbf{By the end of this chapter:} \emph{I can talk about clauses, allocations and following a method.}

Complete the last lesson of chapter 392: I can talk about clauses, allocations and following a method.
\end{chapteropening}

\section[\texorpdfstring{\mr{कलम} (kalam)}{कलम (kalam)}]{\mr{कलम} (kalam) — a clause; a section (of a law)}

\label{lesson:MR-C392-kalam}



\begin{quote}
Before the new one: say the Marathi for a proportion, then the Marathi for a level.
\end{quote}

\subsection*{You'll want to know: \mr{कलम}}
\textbf{\mr{कलम}} — \emph{kalam} — \textquotedblleft{}a clause; a section (of a law)\textquotedblright{}.

\begin{quote}
\mr{करारातील} \mr{हे} \mr{कलम} \mr{मला} \mr{मान्य} \mr{नाही}.
\end{quote}

\begin{quote}
\emph{karārātīl he kalam malā mānya nāhī.}
\end{quote}

\begin{quote}
I do not accept this clause in the contract.
\end{quote}

Neuter.

Neuter: \mr{हे} \mr{कलम}.

\begin{cousinweb}
\textbf{\mr{कलम}} (\emph{kalam}) is a loan from Persian, usually traced to Arabic \emph{qalam}, \textquotedblleft{}pen\textquotedblright{}: an item set down in writing.
\end{cousinweb}

\begin{grammarlens}[title={What you've built}]
One more word for saying what happens if.
\end{grammarlens}

\subsection*{Your turn}
\begin{itemize}
\raggedright
  \item \emph{Say it:} \emph{kalam}
  \item \emph{Say it:} \emph{kalam}, once more
  \item \emph{Say it:} \emph{karārātīl he kalam malā mānya nāhī.} — I do not accept this clause in the contract.
  \item \emph{Recall:} say the Marathi for a proportion, then the Marathi for a level, then say \emph{kalam} again
\end{itemize}

\begin{tcolorbox}[breakable,colback=teal!4,colframe=teal!35!black,title={Before you move on}]
What does \mr{कलम} mean? (\textquotedblleft{}A clause; a section (of a law)\textquotedblright{}.) Say it once more.
\end{tcolorbox}
\section[\texorpdfstring{\mr{वाटप} \mr{करणे} (vāṭap karṇe)}{वाटप करणे (vāṭap karṇe)}]{\mr{वाटप} \mr{करणे} (vāṭap karṇe) — to allocate; to distribute}

\label{lesson:MR-C392-vatap-karne}



\begin{quote}
Before the new one: say the Marathi for a level, then the Marathi for a clause.
\end{quote}

\subsection*{You'll want to know: \mr{वाटप} \mr{करणे}}
\textbf{\mr{वाटप} \mr{करणे}} — \emph{vāṭap karṇe} — \textquotedblleft{}to allocate; to distribute\textquotedblright{}.

\begin{quote}
\mr{संस्थेने} \mr{पूरग्रस्तांना} \mr{अन्नाचं} \mr{वाटप} \mr{केलं}.
\end{quote}

\begin{quote}
\emph{saṁsthene pūragrastānnā annācaṁ vāṭap kelaṁ.}
\end{quote}

\begin{quote}
The organization distributed food to the flood victims.
\end{quote}

Past, with -\mr{ने} (he …): \mr{त्याने} \mr{वाटप} \mr{केलं}, \emph{tyāne vāṭap kelaṁ}.

What is shared out takes -\mr{चं}: \mr{अन्नाचं} \mr{वाटप}.

\begin{cousinweb}
\textbf{\mr{वाटप} \mr{करणे}} (\emph{vāṭap karṇe}) joins \mr{वाटप}, \textquotedblleft{}a sharing out\textquotedblright{}, built on the verb \mr{वाटणे}, \textquotedblleft{}to share out\textquotedblright{} (CDIAL 11238 lists Marathi \emph{vā̃ṭṇẽ} under Sanskrit \emph{vaṇṭati}, \textquotedblleft{}shares\textquotedblright{}), to \mr{करणे}.
\end{cousinweb}

\begin{grammarlens}[title={What you've built}]
One more word for saying what happens if.
\end{grammarlens}

\subsection*{Your turn}
\begin{itemize}
\raggedright
  \item \emph{Say it:} \emph{vāṭap karṇe}
  \item \emph{Say it:} \emph{vāṭap karṇe}, once more
  \item \emph{Say it:} \emph{saṁsthene pūragrastānnā annācaṁ vāṭap kelaṁ.} — The organization distributed food to the flood victims.
  \item \emph{Recall:} say the Marathi for a level, then the Marathi for a clause, then say \emph{vāṭap karṇe} again
\end{itemize}

\begin{tcolorbox}[breakable,colback=teal!4,colframe=teal!35!black,title={Before you move on}]
What does \mr{वाटप} \mr{करणे} mean? (\textquotedblleft{}To allocate; to distribute\textquotedblright{}.) Say it once more.
\end{tcolorbox}
\section[\texorpdfstring{\mr{डावलणे} (ḍāvalṇe)}{डावलणे (ḍāvalṇe)}]{\mr{डावलणे} (ḍāvalṇe) — to pass over; to sideline}

\label{lesson:MR-C392-davalne}



\begin{quote}
Before the new one: say the Marathi for a clause, then the Marathi for to allocate.
\end{quote}

\subsection*{You'll want to know: \mr{डावलणे}}
\textbf{\mr{डावलणे}} — \emph{ḍāvalṇe} — \textquotedblleft{}to pass over; to sideline\textquotedblright{}.

\begin{quote}
\mr{निवड} \mr{करताना} \mr{त्यांनी} \mr{अनुभवी} \mr{लोकांना} \mr{डावललं}.
\end{quote}

\begin{quote}
\emph{nivaḍ kartānā tyānnī anubhvī lokānnā ḍāvallaṁ.}
\end{quote}

\begin{quote}
While choosing, they passed over experienced people.
\end{quote}

Past, with -\mr{ने} (he …): \mr{त्याने} \mr{डावललं}, \emph{tyāne ḍāvallaṁ}.

The one passed over takes -\mr{ला}.

\begin{cousinweb}
\textbf{\mr{डावलणे}} (\emph{ḍāvalṇe}) is a Marathi verb whose origin is not settled.
\end{cousinweb}

\begin{grammarlens}[title={What you've built}]
One more word for saying what happens if.
\end{grammarlens}

\subsection*{Your turn}
\begin{itemize}
\raggedright
  \item \emph{Say it:} \emph{ḍāvalṇe}
  \item \emph{Say it:} \emph{ḍāvalṇe}, once more
  \item \emph{Say it:} \emph{nivaḍ kartānā tyānnī anubhvī lokānnā ḍāvallaṁ.} — While choosing, they passed over experienced people.
  \item \emph{Recall:} say the Marathi for a clause, then the Marathi for to allocate, then say \emph{ḍāvalṇe} again
\end{itemize}

\begin{tcolorbox}[breakable,colback=teal!4,colframe=teal!35!black,title={Before you move on}]
What does \mr{डावलणे} mean? (\textquotedblleft{}To pass over; to sideline\textquotedblright{}.) Say it once more.
\end{tcolorbox}
\section[\texorpdfstring{\mr{अनुसरणे} (anusarṇe)}{अनुसरणे (anusarṇe)}]{\mr{अनुसरणे} (anusarṇe) — to follow (a method)}

\label{lesson:MR-C392-anusarne}



\begin{quote}
Before the new one: say the Marathi for to allocate, then the Marathi for to pass over.
\end{quote}

\subsection*{You'll want to know: \mr{अनुसरणे}}
\textbf{\mr{अनुसरणे}} — \emph{anusarṇe} — \textquotedblleft{}to follow (a method)\textquotedblright{}.

\begin{quote}
\mr{आम्ही} \mr{शिक्षकांनी} \mr{सांगितलेली} \mr{पद्धत} \mr{अनुसरली}.
\end{quote}

\begin{quote}
\emph{āmhī śikṣakānnī sāṅgitlelī paddhat anusarlī.}
\end{quote}

\begin{quote}
We followed the method the teacher described.
\end{quote}

Past, with -\mr{ने} (he …): \mr{त्याने} \mr{अनुसरलं}, \emph{tyāne anusarlaṁ}.

Transitive: the past agrees with what is followed, \mr{पद्धत} \mr{अनुसरली}.

\begin{cousinweb}
\textbf{\mr{अनुसरणे}} (\emph{anusarṇe}) is a learned verb built on Sanskrit \emph{anusarati}, \textquotedblleft{}runs to, runs after\textquotedblright{} (CDIAL 338), with the Marathi ending -\mr{णे}.
\end{cousinweb}

\begin{grammarlens}[title={What you've built}]
One more word for saying what happens if.
\end{grammarlens}

\subsection*{Your turn}
\begin{itemize}
\raggedright
  \item \emph{Say it:} \emph{anusarṇe}
  \item \emph{Say it:} \emph{anusarṇe}, once more
  \item \emph{Say it:} \emph{āmhī śikṣakānnī sāṅgitlelī paddhat anusarlī.} — We followed the method the teacher described.
  \item \emph{Recall:} say the Marathi for to allocate, then the Marathi for to pass over, then say \emph{anusarṇe} again
\end{itemize}

\begin{tcolorbox}[breakable,colback=teal!4,colframe=teal!35!black,title={Before you move on}]
What does \mr{अनुसरणे} mean? (\textquotedblleft{}To follow (a method)\textquotedblright{}.) Say it once more.
\end{tcolorbox}
\section[\texorpdfstring{\mr{अंगीकारणे} (aṅgīkārṇe)}{अंगीकारणे (aṅgīkārṇe)}]{\mr{अंगीकारणे} (aṅgīkārṇe) — to embrace; to adopt}

\label{lesson:MR-C392-angikarne}



\begin{quote}
Before the new one: say the Marathi for to pass over, then the Marathi for to follow.
\end{quote}

\subsection*{You'll want to know: \mr{अंगीकारणे}}
\textbf{\mr{अंगीकारणे}} — \emph{aṅgīkārṇe} — \textquotedblleft{}to embrace; to adopt\textquotedblright{}.

\begin{quote}
\mr{गावाने} \mr{स्वच्छतेचा} \mr{नवा} \mr{मार्ग} \mr{अंगीकारला}.
\end{quote}

\begin{quote}
\emph{gāvāne svacchatecā navā mārga aṅgīkārlā.}
\end{quote}

\begin{quote}
The village embraced a new way of keeping clean.
\end{quote}

Past, with -\mr{ने} (he …): \mr{त्याने} \mr{अंगीकारलं}, \emph{tyāne aṅgīkārlaṁ}.

Formal; \mr{अंगीकार}, adoption.

\begin{cousinweb}
\textbf{\mr{अंगीकारणे}} (\emph{aṅgīkārṇe}) is built on the tatsama \mr{अंगीकार}, Sanskrit \emph{aṅgīkāra}, \textquotedblleft{}agreement, promise\textquotedblright{} (Monier-Williams), built on \emph{aṅga}, \textquotedblleft{}limb, body\textquotedblright{}, with the verb ending -\mr{णे}.
\end{cousinweb}

\begin{grammarlens}[title={What you've built}]
One more word for saying what happens if.
\end{grammarlens}

\subsection*{Your turn}
\begin{itemize}
\raggedright
  \item \emph{Say it:} \emph{aṅgīkārṇe}
  \item \emph{Say it:} \emph{aṅgīkārṇe}, once more
  \item \emph{Say it:} \emph{gāvāne svacchatecā navā mārga aṅgīkārlā.} — The village embraced a new way of keeping clean.
  \item \emph{Recall:} say the Marathi for to pass over, then the Marathi for to follow, then say \emph{aṅgīkārṇe} again
\end{itemize}

\begin{tcolorbox}[breakable,colback=teal!4,colframe=teal!35!black,title={Before you move on}]
What does \mr{अंगीकारणे} mean? (\textquotedblleft{}To embrace; to adopt\textquotedblright{}.) Say it once more.
\end{tcolorbox}
